% Ported verbatim from paper/top3_joint_moment_bound.tex lines 579-620.
% Only cross-reference wording was adapted; no mathematics was altered.
By the preceding section, write $S=\{v_1,\ldots,v_m\}$ with actual masses
$p_i>0$. Put $V=(v_i^T)_{i=1}^m$, $P_0=\diag(p_i)$ and
$B_{ij}=\mathbf1_{\{v_i\cdot v_j>0\}}$. Then
\begin{equation}\label{finite-rows}
 BP_0V=VL,\qquad \sum_i p_iv_iv_i^T=I_3,\qquad
 v_i^TLv_i=\mathcal M.
\end{equation}
The full rows of $B$ are distinct: equal rows would imply equal $v_i$
by invertibility of $L$. A positive clique $J$ means a subset of this actual support with $B_{ij}=1$
for every $i,j\in J$. Extend vectors indexed by $J$ by zero outside $J$.
\begin{lemma}[Vector rigidity on a positive clique]
For every real $a$ supported on a positive clique, if
$Q_a=\sum_i a_i v_iv_i^T=0$, then
\begin{equation}
 \sum_i a_i=0,\quad b_a:=\sum_i a_iv_i=0,\quad
 \sum_i B_{ki}a_iv_i=0\quad\hbox{for every row }k. \label{rigid}
\end{equation}
\end{lemma}
\begin{proof}
Taking the trace against $L$ gives $\mathcal M\sum a_i=0$.
For both signs of sufficiently small nonzero $t$, all $p_i+ta_i$ are positive.
These are legitimate new fiber measures for the same binary template, realizable
on an atomless space. They are not independent choices of moment data.
The original coordinate functions have Gram matrix $I+tQ_a=I$ and compression
\[
 N(t)=L+t(Q_aL+LQ_a)+t^2b_ab_a^T=L+t^2b_ab_a^T.
\]
The last term uses the clique assumption, and both cross terms use \eqref{finite-rows}.
Their trace cannot exceed the global maximum $\mathcal M$; hence $b_a=0$ and $N(t)=L$.
Equality of this trial trace with the maximum forces its three-dimensional
space to be invariant. Explicitly, rotate one orthonormal coordinate into any
unit perpendicular function; the first derivative of the trace gives zero
off-space matrix element. This argument does not differentiate ordered eigenvalues.
The compression $L$ then gives the eigenfunction equations at the perturbed measure.
All fibers still have positive mass, so for every $k$,
$\sum_i(p_i+ta_i)B_{ki}v_i=Lv_k$.
Subtract \eqref{finite-rows} and divide by $t$ to obtain \eqref{rigid}.
\end{proof}
Moment-preserving cloning and the Ky Fan equality argument follow the
mechanism in~\cite[Claims 2.2--2.5]{KMMP}. The permissible size of $t$
may depend on the maximizer and the chosen direction.
